\section{Data Sources and Collection}
\label{sec:data}

\subsection{Sources and licences}

The project combines one domain-specific API with two general LOD datasets.
Table~\ref{tab:sources} lists what each source contributes, how it is accessed
and under which licence. Only the documented API and SPARQL endpoints are used;
no review sites are scraped. DBpedia is not queried during collection: its
resource URIs are derived later from the English Wikipedia sitelinks of matched
Wikidata items (Section~\ref{sec:interlinking}).

\begin{table}[h]
  \centering
  \small
  \begin{tabular}{@{}p{2.5cm}p{5.6cm}p{2.6cm}p{3.3cm}@{}}
    \toprule
    Source & Content used & Access & Licence \\
    \midrule
    TheCocktailDB \cite{cocktaildb} & Cocktail names, IBA flag, glass,
    instructions, images, ingredient lists and measures (441 drinks) & Public
    JSON API & Free for non-commercial use, attribution \\
    Wikidata \cite{wikidata} & Distilleries (\texttt{P31 Q1251750}) with
    coordinates, country, inception, website, image and enwiki sitelink;
    spirit brands (\texttt{P31/P279* Q56139} with manufacturer a distillery)
    & SPARQL endpoint (collection) and MediaWiki API (linking) & CC0 1.0 \\
    DBpedia \cite{dbpedia} & Resource URIs used as \texttt{owl:sameAs} targets
    & Resource URIs (no bulk download) & CC BY-SA 4.0 \\
    \bottomrule
  \end{tabular}
  \caption{Data sources, content and licences; full details are kept in
  \texttt{docs/data-licences.md}.}
  \label{tab:sources}
\end{table}

\subsection{The collection pipeline}

Figure~\ref{fig:pipeline} shows the end-to-end pipeline. The collection step is
implemented in \texttt{etl/bootstrap.py}. Because TheCocktailDB offers no ``dump
everything'' endpoint, the collector queries the per-letter search endpoint for
\texttt{a}--\texttt{z} and \texttt{0}--\texttt{9} and deduplicates the results by
the source identifier \texttt{idDrink}, yielding 441 cocktails. Wikidata is
queried directly in SPARQL: the distillery query selects items whose
\texttt{P31} is \texttt{Q1251750} with optional coordinates (\texttt{P625}),
country (\texttt{P17}), inception (\texttt{P571}), website (\texttt{P856}),
image (\texttt{P18}) and English article, while the brand query uses the
property path \texttt{P31/P279*} to include spirit subclasses and requires the
manufacturer (\texttt{P176}) to be a distillery.

\begin{figure}[h]
  \centering
  \begin{tikzpicture}[
    node distance=4.2mm,
    box/.style={draw, rounded corners=2pt, fill=blue!5, align=center,
                font=\small, inner sep=3pt, minimum height=7.5mm},
    src/.style={draw, rounded corners=2pt, fill=orange!10, align=center,
                font=\small, inner sep=3pt, minimum height=7.5mm},
    arr/.style={-{Latex[length=1.8mm]}, thick},
    lbl/.style={font=\scriptsize\itshape, right=2mm, align=left}
  ]
    \node[src] (cdb) {TheCocktailDB JSON API};
    \node[src, right=6mm of cdb] (wd) {Wikidata SPARQL endpoint};
    \node[box, below=6mm of $(cdb.east)!0.5!(wd.west)$] (boot)
      {\texttt{etl/bootstrap.py}: collect, cache, write manifest};
    \node[box, below=4.2mm of boot] (raw)
      {\texttt{data/raw/} JSON dumps $+$ \texttt{manifest.json}};
    \node[box, below=4.2mm of raw] (map)
      {\texttt{etl/map.py}: transform JSON to RDF};
    \node[box, below=4.2mm of map] (data)
      {\texttt{data/rdf/data.rdf} (instance graph, ontology terms)};
    \node[box, below=4.2mm of data] (rec)
      {\texttt{etl/reconcile.py}: match against Wikidata and DBpedia};
    \node[box, below=4.2mm of rec] (links) {\texttt{data/links/links.rdf}};
    \node[box, below=4.2mm of links] (serve)
      {\texttt{etl/run\_queries.py} $\cdot$ \texttt{web/server.py} $\cdot$ Fuseki};
    \draw[arr] (cdb) -- (boot);
    \draw[arr] (wd) -- (boot);
    \draw[arr] (boot) -- (raw);
    \draw[arr] (raw) -- (map);
    \draw[arr] (map) -- (data);
    \draw[arr] (data) -- (rec);
    \draw[arr] (rec) -- (links);
    \draw[arr] (links) -- (serve);
  \end{tikzpicture}
  \caption{The end-to-end pipeline: collection, transformation, interlinking
  and serving. Earlier stages are cached, so later stages can be re-run
  without network access.}
  \label{fig:pipeline}
\end{figure}

\subsection{Designing the Wikidata queries}
\label{sec:wikidata-queries}

The two queries were derived from the ontology rather than written ad hoc:
every class and field of Section~\ref{sec:ontology} dictates either a class
constraint or an optional value to fetch. Class constraints use \texttt{P31}
(\emph{instance of}): a distillery is an item whose \texttt{P31} is
\texttt{Q1251750}. Spirit brands need one extra step, because a brand is an
instance of a \emph{specific} spirit class (gin, whisky, \dots) and those
classes are \texttt{P279} (\emph{subclass of}) \texttt{Q56139} (liquor). The
property path \texttt{wdt:P31/wdt:P279*} therefore reads ``instance of liquor
or of any of its subclasses'', where \texttt{*} allows zero or more subclass
steps; the path was verified on Wikidata (gin, \texttt{Q170210}, is a subclass
of \texttt{Q56139}). The manufacturer constraint \texttt{?b wdt:P176 ?d} plus
\texttt{?d wdt:P31 wd:Q1251750} is what ties a brand to a distillery and makes
the \texttt{drink:producedBy} edge of Section~\ref{sec:ontology} possible.

All descriptive fields are wrapped in \texttt{OPTIONAL} because Wikidata is
incomplete: a missing coordinate or website must not remove the entity from the
result. Readable labels are produced by the standard
\texttt{SERVICE wikibase:label} block, which turns each \texttt{?x} variable
into an \texttt{?xLabel} from \texttt{rdfs:label} using a language fallback
list. The English Wikipedia article is fetched with the \texttt{schema:about}
plus \texttt{schema:isPartOf} pattern, because sitelinks are not ordinary
statements; the resulting URL is what Section~\ref{sec:interlinking} converts
into a DBpedia URI. Listing~\ref{lst:wikidata-queries} shows both queries as
executed by the collector, and Table~\ref{tab:wikidata-props} summarises the
properties they use.

\begin{lstlisting}[frame=single, basicstyle=\ttfamily\scriptsize, breaklines=true,
    caption={The two Wikidata collection queries: spirit brands (top) and
    distilleries (bottom).},
    label={lst:wikidata-queries}]
SELECT ?b ?bLabel ?d ?dLabel ?country ?countryLabel ?inception ?website ?image ?article WHERE {
  ?b wdt:P31/wdt:P279* wd:Q56139 .
  ?b wdt:P176 ?d .
  ?d wdt:P31 wd:Q1251750 .
  OPTIONAL { ?b wdt:P495 ?country . }
  OPTIONAL { ?b wdt:P571 ?inception . }
  OPTIONAL { ?b wdt:P856 ?website . }
  OPTIONAL { ?b wdt:P18 ?image . }
  OPTIONAL { ?article schema:about ?b ; schema:isPartOf <https://en.wikipedia.org/> . }
  SERVICE wikibase:label { bd:serviceParam wikibase:language "en". }
}

SELECT ?d ?dLabel ?coord ?country ?countryLabel ?inception ?website ?image ?article WHERE {
  ?d wdt:P31 wd:Q1251750 .
  OPTIONAL { ?d wdt:P625 ?coord . }
  OPTIONAL { ?d wdt:P17 ?country . }
  OPTIONAL { ?d wdt:P571 ?inception . }
  OPTIONAL { ?d wdt:P856 ?website . }
  OPTIONAL { ?d wdt:P18 ?image . }
  OPTIONAL { ?article schema:about ?d ; schema:isPartOf <https://en.wikipedia.org/> . }
  SERVICE wikibase:label { bd:serviceParam wikibase:language "en,de,fr,es". }
}
\end{lstlisting}

\begin{table}[h]
  \centering
  \small
  \begin{tabular}{@{}lll@{}}
    \toprule
    Property & Meaning & Used in \\
    \midrule
    \texttt{P31} & instance of & both queries \\
    \texttt{P279} & subclass of & brand query (property path) \\
    \texttt{P176} & manufacturer & brand query \\
    \texttt{P625} & coordinate location & distillery query \\
    \texttt{P17} & country & distillery query \\
    \texttt{P495} & country of origin & brand query \\
    \texttt{P571} & inception & both queries \\
    \texttt{P856} & official website & both queries \\
    \texttt{P18} & image & both queries \\
    \bottomrule
  \end{tabular}
  \caption{Wikidata properties used by the two queries. The class constraints
  are \texttt{Q1251750} (distillery) and \texttt{Q56139} (liquor).}
  \label{tab:wikidata-props}
\end{table}

\subsection{Sample query results}
\label{sec:query-samples}

Tables~\ref{tab:sample-distilleries} and \ref{tab:sample-brands} show
representative rows of the two queries as mapped into the graph: the distillery
query returned 998 rows for 865 distinct items, and the brand query returned 12
rows. Optional fields are frequently absent in Wikidata, which is why several
cells in the tables are empty.

\begin{table}[h]
  \centering
  \small
  \begin{tabular}{@{}lllr@{}}
    \toprule
    Distillery & Coordinates (lat, lon) & Country & Founded \\
    \midrule
    Highland Park distillery & 58.969, -2.955 & United Kingdom & 1788 \\
    Frapin & 45.617, -0.218 & France & 1270 \\
    John Watling's Distillery & 25.075, -77.348 & The Bahamas & 1789 \\
    Loch Ewe Distillery & 57.834, -5.574 & United Kingdom & 2004 \\
    \bottomrule
  \end{tabular}
  \caption{Sample rows from the distillery query (\texttt{?dLabel},
  \texttt{P625}, country label, \texttt{P571}).}
  \label{tab:sample-distilleries}
\end{table}

\begin{table}[h]
  \centering
  \small
  \begin{tabular}{@{}lll@{}}
    \toprule
    Brand & Manufacturer (distillery) & Website (\texttt{P856}) \\
    \midrule
    Gentleman Jack & Jack Daniel's & n/a \\
    Kyr\"o Napue Gin & Kyr\"o Distillery Company & n/a \\
    Connemara & Cooley Distillery & n/a \\
    Viru Valge & Liviko & viruvalge.ee \\
    \bottomrule
  \end{tabular}
  \caption{Sample rows from the brand query: the brand label, its manufacturer
  (the \texttt{?dLabel} that becomes \texttt{drink:producedBy}) and the
  optional website field, shown as host only.}
  \label{tab:sample-brands}
\end{table}

\subsection{Provenance and reproducibility}

Every download is cached under \texttt{data/raw/} and skipped on re-runs unless
\texttt{--force} is given; the pipeline is therefore idempotent. For each dump,
the collector writes an entry to \texttt{data/raw/manifest.json} containing the
source (a URL or a query description), licence, retrieval timestamp, record
count and a SHA-256 hash of the file. Table~\ref{tab:manifest} summarises the three dumps. The hash makes
changes detectable: re-downloading a source that has changed produces a
different fingerprint, which is exactly the kind of provenance evidence the
first four stars require.

\begin{table}[h]
  \centering
  \small
  \begin{tabular}{@{}llrl@{}}
    \toprule
    Dump & Source & Records & SHA-256 (first 12 hex) \\
    \midrule
    \texttt{cocktaildb/drinks.json} & TheCocktailDB per-letter search & 441 & \texttt{1ba6721fc45d} \\
    \texttt{wikidata/distilleries.json} & Wikidata SPARQL (\texttt{Q1251750}) & 998 & \texttt{53fb92019434} \\
    \texttt{wikidata/brands.json} & Wikidata SPARQL (spirit brands) & 12 & \texttt{e483ee7ee53c} \\
    \bottomrule
  \end{tabular}
  \caption{Contents of the provenance manifest after collection.}
  \label{tab:manifest}
\end{table}

\subsection{Verifying external identifiers}

The project plan initially identified distilleries by \texttt{Q18979992}. Before
implementing the collector, the identifier was checked against Wikidata and
turned out to denote \emph{a street in the Netherlands}. The correct class is
\texttt{Q1251750}. This episode shaped a habit that appears again in
Section~\ref{sec:interlinking}: external identifiers are verified against the source, never trusted
from documentation alone.

\subsection{Characteristics of the raw data}

Several properties of the sources influenced the later design. The distillery
SPARQL result contains 998 rows but only 865 distinct distilleries, because
items with several values (for example multiple coordinates) produce multiple
bindings; the transformation step therefore groups rows by QID first. Optional
fields are frequently missing: some cocktails have no measure for an
ingredient, and measures are free text such as ``Juice of 1/2'' rather than
numbers, which is why Section~\ref{sec:transformation} keeps the original text while parsing a numeric
value only when possible. Finally, only 12 spirit brands are returned, because
the requirement that the manufacturer be recorded as a distillery is rarely
satisfied in Wikidata; this limitation is revisited in Section~\ref{sec:evaluation}.
